% ===================================================================== PREAMBULE COMMUN — Carnet d'entraînement Fiche 9
\documentclass[11pt,a4paper]{article}
\usepackage[utf8]{inputenc}
\usepackage[T1]{fontenc}
\usepackage{lmodern}
\usepackage[french]{babel}
\usepackage{amsmath,amssymb,mathtools}
\usepackage{icomma}
\usepackage{siunitx}
\sisetup{output-decimal-marker={,},per-mode=symbol}
\DeclareSIUnit{\molar}{M}
\DeclareSIUnit{\Molar}{mol\,L^{-1}}
\usepackage[version=4]{mhchem}
\usepackage{xcolor}
\usepackage{colortbl}
\usepackage{tikz}
\usepackage{pgfplots}
\usepackage[most]{tcolorbox}
\usepackage{enumitem}
\usepackage{array}
\usepackage{booktabs}
\usepackage{multicol}
\usepackage{fancyhdr}
\usepackage[hidelinks]{hyperref}
\usetikzlibrary{arrows.meta,positioning,calc,decorations.pathmorphing,
  decorations.markings,angles,quotes,patterns,backgrounds,
  shapes.geometric,shapes.misc,fadings,intersections,fit}
\pgfplotsset{compat=1.18}
\usepackage[a4paper,margin=2.2cm,headheight=26pt,headsep=10pt]{geometry}

% ---------- Palette ----------
\definecolor{primary}{RGB}{150,98,162}
\definecolor{primarydark}{RGB}{92,52,110}
\definecolor{defcol}{RGB}{0,114,178}
\definecolor{thmcol}{RGB}{0,140,120}
\definecolor{formulacol}{RGB}{198,93,9}
\definecolor{remarkcol}{RGB}{120,94,200}
\definecolor{attncol}{RGB}{200,40,55}
\definecolor{excol}{RGB}{0,135,90}
\definecolor{methcol}{RGB}{90,90,100}
\definecolor{cationcol}{RGB}{0,90,190}
\definecolor{anioncol}{RGB}{200,60,55}
\definecolor{cristalcol}{RGB}{110,105,120}
\definecolor{eaucol}{RGB}{120,180,220}
\definecolor{solcol}{RGB}{0,135,90}
\definecolor{kpscol}{RGB}{198,93,9}
\definecolor{qcol}{RGB}{120,94,200}
\definecolor{precipcol}{RGB}{110,70,40}
\definecolor{exercol}{RGB}{150,98,162}
\definecolor{corrcol}{RGB}{0,120,100}

% ---------- Macros ----------
\newcommand{\Kps}{K_{\mathrm{ps}}}
\newcommand{\Ce}[1]{\left[\,\ce{#1}\,\right]}

% ---------- Style des boîtes ----------
\tcbset{boxbase/.style={enhanced,breakable,boxrule=0.9pt,arc=2.5pt,
   left=3mm,right=3mm,top=2.3mm,bottom=2.3mm,
   fonttitle=\bfseries,coltitle=white,toptitle=0.6mm,bottomtitle=0.6mm}}

\newtcolorbox[auto counter]{definition}[1][]{%
  boxbase,colback=defcol!5,colframe=defcol,
  title={\textsc{Définition}~\thetcbcounter\ \ifx\relax#1\relax\relax\else\textmd{--- \itshape #1}\fi}}
\newtcolorbox[auto counter]{theoreme}[1][]{%
  boxbase,colback=thmcol!6,colframe=thmcol,
  title={\textsc{Loi / Propriété}~\thetcbcounter\ \ifx\relax#1\relax\relax\else\textmd{--- \itshape #1}\fi}}
\newtcolorbox[auto counter]{formule}[1][]{%
  boxbase,colback=formulacol!6,colframe=formulacol,
  title={\textsc{Formule}~\thetcbcounter\ \ifx\relax#1\relax\relax\else\textmd{--- \itshape #1}\fi}}
\newtcolorbox{remarque}[1][]{%
  boxbase,colback=remarkcol!6,colframe=remarkcol,
  title={\textsc{Remarque}\ifx\relax#1\relax\relax\else\textmd{ : \itshape #1}\fi}}
\newtcolorbox{attention}[1][]{%
  boxbase,colback=attncol!6,colframe=attncol,
  title={\textsc{Attention}\ifx\relax#1\relax\relax\else\textmd{ : \itshape #1}\fi}}
\newtcolorbox{exemple}[1][]{%
  boxbase,colback=excol!5,colframe=excol,
  title={\textsc{Exemple}\ifx\relax#1\relax\relax\else\textmd{ : \itshape #1}\fi}}
\newtcolorbox{methode}[1][]{%
  boxbase,colback=methcol!7,colframe=methcol,sharp corners,
  title={\textsc{Méthode}\ifx\relax#1\relax\relax\else\textmd{ : \itshape #1}\fi}}
\newtcolorbox{astuce}[1][]{%
  boxbase,colback=primary!7,colframe=primary,
  title={\textsc{Astuce}\ifx\relax#1\relax\relax\else\textmd{ : \itshape #1}\fi}}

\newtcolorbox[auto counter]{exercice}[1][]{%
  boxbase,colback=exercol!4,colframe=exercol,
  title={\textsc{Exercice}~\thetcbcounter\ \ifx\relax#1\relax\relax\else\textmd{--- \itshape #1}\fi}}
\newtcolorbox[auto counter]{correction}[1][]{%
  boxbase,colback=corrcol!4,colframe=corrcol,
  title={\textsc{Corrigé}~\thetcbcounter\ \ifx\relax#1\relax\relax\else\textmd{--- \itshape #1}\fi}}

% ---------- Environnement « où : » ----------
\newenvironment{ou}{%
  \par\smallskip\noindent\textbf{où :}\nopagebreak
  \begin{itemize}[leftmargin=1.8em,topsep=2pt,itemsep=1.5pt,parsep=0pt]}%
  {\end{itemize}}

% ---------- Styles TikZ ----------
\tikzset{
  cat/.style={circle,fill=cationcol,draw=cationcol!50!black,inner sep=0pt,
              minimum size=5mm,font=\bfseries\scriptsize,text=white},
  an/.style={circle,fill=anioncol,draw=anioncol!50!black,inner sep=0pt,
             minimum size=5mm,font=\bfseries\scriptsize,text=white},
  crx/.style={circle,fill=cristalcol,draw=black!50,inner sep=0pt,
              minimum size=5mm,font=\bfseries\scriptsize,text=white},
  ptn/.style={circle,fill=black,inner sep=1.1pt}}

% ---------- En-têtes ----------
\pagestyle{fancy}
\fancyhf{}
\fancyhead[L]{\small\textcolor{primarydark}{\textbf{Fiche 9} — Solubilité et produit de solubilité}}
\fancyhead[R]{\small\textcolor{primarydark}{\textsc{Carnet d'entraînement}}}
\fancyfoot[C]{\small\thepage}
\renewcommand{\headrulewidth}{0.8pt}
\renewcommand{\headrule}{\hbox to\headwidth{\color{primary}\leaders\hrule height \headrulewidth\hfill}}
\usepackage{titlesec}
\titleformat{\section}{\Large\bfseries\color{primarydark}}{\thesection}{0.6em}{}
\titleformat{\subsection}{\large\bfseries\color{primary}}{\thesubsection}{0.6em}{}

% ---------- Bandeau de titre ----------
% #1 = thème/étiquette   #2 = titre principal   #3 = sous-titre
\newcommand{\bandeau}[3]{%
\begin{center}
\begin{tikzpicture}
\node[rectangle,rounded corners=7pt,fill=primary,text=white,
      minimum width=15.6cm,minimum height=1.95cm,align=center,inner sep=8pt]
  {{\large\bfseries #1}\\[3pt]{\Large\bfseries #2}\\[3pt]{\small #3}};
\end{tikzpicture}
\end{center}
\vspace{2mm}}
\begin{document}
\bandeau{Thème 2 — Relation $\Kps \leftrightarrow s$}%
{Corrigé — Niveau Difficile}%
{Réponses détaillées et justifications}

\begin{correction}[dérivation complète — \ce{Ca3(PO4)2}]
\begin{enumerate}[label=\textbf{\alph*)},leftmargin=2.2em,topsep=2pt,itemsep=2pt]
  \item $\ce{Ca3(PO4)2(s) <=> 3 Ca^{2+}(aq) + 2 PO4^{3-}(aq)}$.
  \item À l'équilibre : $[\ce{Ca^2+}]=\mathbf{3s}$ et $[\ce{PO4^3-}]=\mathbf{2s}$.
  \item $\Kps=[\ce{Ca^2+}]^{3}[\ce{PO4^3-}]^{2}=(3s)^{3}(2s)^{2}=27s^{3}\cdot 4s^{2}=\boxed{108\,s^{5}}$.
  \item $s=\sqrt[5]{\dfrac{\Kps}{108}}=\sqrt[5]{\dfrac{\num{1{,}3e-32}}{108}}
        =\sqrt[5]{\num{1{,}2e-34}}\approx\boxed{\qty{1{,}6e-7}{\molar}}$.\\
        Puis $[\ce{PO4^3-}]=2s\approx\qty{3{,}3e-7}{\molar}$ (et $[\ce{Ca^2+}]=3s\approx\qty{4{,}9e-7}{\molar}$).\\
        \emph{Contrôle :} $108\,(1{,}6\times10^{-7})^{5}\approx\num{1{,}1e-32}$, du bon ordre de grandeur.
\end{enumerate}
\end{correction}

\begin{correction}[oxalate de calcium]
\begin{enumerate}[label=\textbf{\alph*)},leftmargin=2.2em,topsep=2pt,itemsep=2pt]
  \item $\ce{CaC2O4(s) <=> Ca^{2+}(aq) + C2O4^{2-}(aq)}$ ; type $1{:}1$ donc $\Kps=[\ce{Ca^2+}][\ce{C2O4^2-}]=s^{2}$.
  \item $s=\qty{3{,}0e-6}{\molar}$ signifie que \qty{1}{\litre} d'eau saturée dissout \emph{au maximum}
        \qty{3{,}0e-6}{\mole} de \ce{CaC2O4}.
  \item $n=s\,V\Rightarrow V=\dfrac{n}{s}=\dfrac{\qty{6{,}0e-3}{\mole}}{\qty{3{,}0e-6}{\mole\per\litre}}
        =\boxed{\qty{2000}{\litre}}$ (deux mètres cubes !). \emph{(Réponse C du \textsc{qcm}\,n\textsuperscript{o}\,10.)}
  \item $s_m=s\times M=\num{3{,}0e-6}\times 128\approx\qty{3{,}8e-4}{\gram\per\litre}$, soit \qty{0{,}38}{\milli\gram\per\litre}.
        C'est infime : une fois formé, un calcul ne se dissout pas par simple dilution — d'où le recours
        à la lithotripsie (fragmentation par ultrasons/laser) en médecine.
\end{enumerate}
\end{correction}

\begin{correction}[iodate de strontium \ce{Sr(IO3)2}]
\begin{enumerate}[label=\textbf{\alph*)},leftmargin=2.2em,topsep=2pt,itemsep=2pt]
  \item $\ce{Sr(IO3)2(s) <=> Sr^{2+}(aq) + 2 IO3-(aq)}$. Une unité libère 1 \ce{Sr^{2+}} pour
        \textbf{2} \ce{IO3-} : il y a \emph{deux fois plus} d'ions iodate. L'affirmation d'annales
        « autant d'ions strontium que d'ions iodate » est donc \textbf{fausse}.
  \item $[\ce{Sr^2+}]=s=\qty{9{,}0e-4}{\molar}$ ; $[\ce{IO3-}]=2s=\qty{1{,}8e-3}{\molar}$.
  \item $\Kps=[\ce{Sr^2+}][\ce{IO3-}]^{2}=(s)(2s)^{2}=4\,s^{3}
        =4\,(\num{9{,}0e-4})^{3}=4\times\num{7{,}29e-10}\approx\boxed{\num{2{,}9e-9}}$.
  \item Charge des cations : $2\times[\ce{Sr^2+}]=2\times\num{9{,}0e-4}=\num{1{,}8e-3}$ ;
        charge des anions : $1\times[\ce{IO3-}]=\num{1{,}8e-3}$. Elles sont égales : la solution est bien
        \textbf{électriquement neutre}. (Avoir un rapport $1{:}2$ n'empêche pas la neutralité, car les
        charges — $+2$ contre $-1$ — se compensent.)
\end{enumerate}
\end{correction}

\end{document}
